\chapter{Discussion}\label{ch:discussion}

\section{Power law behaviour}\label{sec:discuss-powerlaw}

Power-law behaviours are theoretically plausible for topples, area and length. This is because an SOC state has no intrinsic characteristic avalanche scale, as we previously discussed in Chapter \ref{ch:background}. The distribution therefore behaves approximately as a power law and is truncated by the finite lattice. Topples measure the total ``activity" generated during the process that take the system back to a stable configuration, area measures the number of distinct sites affected by this avalanche process, and length measures the distance this avalanche travels. Hence if an avalanche remains scale invariant, all these closely related properties should remain scale invariant and are hence expected to exhibit approximately scale-free distributions within the truncated region. This argument does not apply as directly to loss, because loss also depends specifically on whether an avalanche reaches the open boundary, which we will discuss in detail under Section \ref{sec:discuss-exp}.\\

From our simulation results, the observed topples and area distributions are consistent with this general picture. Figures \ref{fig:topples} and \ref{fig:area} illustrates an approximate linear segment on log-log axes, followed by system size dependent tails. This indicates that avalanches occur over a wide range of scales while the largest observed topples and area are restricted by the finite lattice. \\

For the BTW model, the unscaled estimates $\tau_s=1.120$ and $\tau_a=1.106$ are lower than many of the literature estimates at around 1.3. The corresponding Manna estimates $\tau_s=1.123$ and $\tau_a=1.133$ are also lower than previously reported values \cite{Lubeck1997Numerical}. Several methodology differences may contribute to this discrepancy, including our much smaller lattice sizes, the selected fitting cutoff intervals, and differences in exponent-estimation methods. \\

The interpretation of length is more difficult. The unscaled estimates were $\tau_l=0.884$ for BTW and $\tau_l=0.934$ for Manna. But since our maximum Manhattan distance is not equivalent to the radius of gyration studied by L\"{u}beck and Usadel, our $\tau_l$ is not directly comparable with their reported $\tau_r$. Moreover, the visible curvature and the lower threshold ($R^2 \geq 0.88$) required for length indicate weaker evidence for a power-law region. Hence, these estimates may not accurately reflect the asymptotic length exponent.\\

In addition, the Manna topples and area distributions exhibited an initial plateau. This can reflect discreteness or that very small local avalanches follow a different distribution. Excluding $x=1$ produced only a small change in the estimated exponents (e.g., 1.123 to 1.121 for topples and 1.133 to 1.127 for area, before scaling), suggesting limited effect to this particular lower-cutoff choice. The selected upper cutoffs, however, were more sensitive to the chosen $R^2$ threshold and could extend into the sparse finite-size tail. \\

Overall, for topple, area and length, the intermediate power-law like regions and the extension of the cutoffs with $L$ (as shown in Table \ref{tab:powerlaw-cutoffs}) are consistent with scale-invariant avalanche behaviour and hence with SOC. Nevertheless, the evidence could be weaker for length, and exponent estimates for Manna topples and area may have been affected by cutoff selections. Stronger evidence would require larger lattices, wider scaling regions, uncertainty estimates, and evidence that the fitted exponents stabilise as $L$ increases.

\section{Effect of system-size dependent scaling}\label{sec:discuss-scaling}

In this report, we tested system-size dependent scaling proposed by Yadav et al \cite{Yadav2022FiniteSize} to topples, area and length. We wanted to investigate if this will lead to an improvement in collapse. That is, do the distributions for different system sizes lie closer to a common fitted line? We assessed this using two measures, an overall RMSE and maximum absolute deviations across all $L$. They are complementary with each other as RMSE assesses the overall agreement while the max deviations identify which system sizes improves or deteriorates.\\ 

For the BTW topples and area distributions, the $\log N$ adjustment produced modest improvements. The overall RMSE decreased by approximately 3.9\% for topples and 5.2\% for area, with the largest reductions in maximum deviation occurring for $L=10$ and $L=20$. The adjusted exponents were also closer to the values $\tau_s = 1$ and $\tau_a \approx 1.02$ proposed by Yadav et al. \cite{Yadav2022FiniteSize}. Our findings are consistent with a logarithmic system-size dependence, although the small RMSE reductions and limited range of $L$ cannot provide conclusive evidence.\\

The similar behaviour of topples and area is reasonable because avalanches affecting more distinct sites generally generate more toppling events. Their close relationship is also reflected by the near-perfect Spearman correlations observed for both models.\\

We additionally applied the $\log N$ adjustment to length, which was not considered in the original analysis by Yadav et al. The result is rather interesting as for BTW, the overall RMSE decreased massively by approximately 32\%, and the maximum deviation decreased for every system size. This was the largest relative improvement among BTW observables. Nevertheless, improved alignment does not by itself establish that length follows a power law, especially because our original length distributions exhibited visible curvature.\\

The Manna results were less uniform. For topples, the overall RMSE increased by approximately 2.7\%, and the maximum deviations showed a mixture of increases and decreases. There is therefore no consistent evidence that the adjustment improves the Manna topples collapse. For area, the RMSE decreased by approximately 1.9\%, with most of the improvement again occurring for the smaller system sizes. This provides weak evidence of improved alignment. \\

For Manna length, the overall RMSE decreased by approximately 25\% and the maximum deviations decreased for every system size. As in the BTW case, this suggests that the adjustment reduces vertical separation between the length distributions. However, because the underlying power-law evidence is weaker, this does not by itself imply power-law behaviour.\\

Overall, the $\log N$ adjustment is most consistently supported by BTW area and topples and by the empirical alignment of length. Its effect on the Manna model can depend on avalanche property, since result is inconclusive for Manna topples, weak supportive evidence for Manna area and improved alignment for Manna length.

\section{Exponential behaviour}\label{sec:discuss-exp}

Unlike topples, area and length, loss only happens when an avalanche reaches an open boundary. It is therefore not as directly related to the scale-free characteristics of an avalanche. \\

The loss results provide partial support for an exponential behaviour. The positive-loss probabilities appear approximately linear on log-normal axes, and the corresponding $R^2$ values are relatively high. Moreover, $P(X=0)$ increases with system size for both models. This is consistent with the decreasing proportion of boundary sites over total sites as $L$ increases, since total sites increases proportionally to $L^2$ while boundary site only increases proportionally to $L$. Therefore, a randomly initiated avalanche must propagate further on average to reach an open boundary. The Manna model has a lower zero-loss probability than the BTW model at every system size, indicating that boundary dissipation occurs more frequently in the simulated Manna avalanches.\\

Nevertheless, the residual plots show a systematic U-shaped pattern. The fitted model generally underestimates the probabilities of small and very large losses while overestimating probabilities in the intermediate range. Therefore, the relatively high $R^2$ values should not be seen as evidence of an exact exponential distribution. A single decay rate $\lambda$ hence cannot fully describe its shape. \\

An exponential model was also tested for avalanche length because the unscaled length distributions showed noticeable curvature and weaker power-law linearity than the topples and area distributions. Although some of the log-normal fits produced moderately high $R^2$ values, their residuals exhibited a clear S-shaped pattern. This systematic curvature shows that a single exponential distribution is not an adequate model for positive length.\\

The failure of the exponential model for length is actually also consistent with its earlier power-law interpretation. Length measures the distance travelled by an avalanche and is therefore more directly connected to scale-invariant avalanche property than loss. This can be supported by decreasing zero-length probabilities and decay rate as $L$ increase, which may indicate a broader range of length values.\\

Therefore, the results favour an approximate power-law description for length within a restricted finite-size scaling region over a single exponential description of the complete distribution. However, the visible curvature and the lower cutoff-selection threshold still mean that the evidence for power-law behaviour in length remains weaker than for area and topples

\section{BTW vs. Manna model}\label{sec:discuss-models}

Despite their different toppling rules, both models displayed several common features. Their topples and area distributions contained approximately power-law intermediate regions followed by finite-size cutoffs, and their fitted exponents were relatively close to each other across the system sizes we considered. Both models also exhibited strong positive correlations among topples, area, and length. \\

Nevertheless, several differences were also observed. The Manna distributions exhibited an initial plateau that was not present or at least as prominent in the BTW results. Notably, the $\log N$ adjustment did not consistently improve the Manna topples collapse. The Manna model also produced a lower probability of zero loss, suggesting that its stochastic toppling rule generated avalanches that reach boundaries more frequently. These differences indicate that the toppling rule can affect of the finite-size distributions for some avalanche properties.

\section{Limitations}\label{sec:limitations}

The main computational limitation of this exploration is the restricted range of system sizes. Simulations were limited to $L=50$, whereas  literature exponent estimates usually used substantially larger lattices. Therefore the fitted slopes may not represent their asymptotic values. Attempts were made to resolve this computational issue. We used a queue-based implementation to replace the initial whole grid-search strategy and improved the toppling efficiency. Nevertheless, the computational cost of the Manna model restricted further increases in $L$ since we wanted our simulation results to be easily reproducible with a relatively short amount of time.\\

The exponent estimates are also sensitive to the selected fitting regions. Upper cutoffs were chosen using an $R^2$ threshold applied to the empirical survival function, but this could allow some selected fitting regions to extend into the sparse finite-size tail. Furthermore, OLS was applied to unbinned log-PMF points, giving each distinct observed value equal weight even when tail values occurred rarely, hence potentially biasing the exponent estimates. Alternative cutoff thresholds or likelihood-based power-law methods could therefore produce different estimates.\\

Finally, each result was based on one reproducible simulated trajectory for each model and system size. Although 50000 post-burn-in observations were retained, repeated simulations with independent seeds would be needed to quantify uncertainty intervals. Our collapse analysis also measured agreement with a common fitted line after vertical $\log N$ rescaling rather than performing a complete finite-size collapse involving estimation of the horizontal scaling dimension $D_x$. In addition, our Manhattan-distance definition of length is not directly comparable with the radius of gyration commonly used in the literature.
